\documentclass[output=paper]{langsci/langscibook}
\ChapterDOI{10.5281/zenodo.4280637}
\author{George Tsoulas\affiliation{University of York}}
\title{On a difference between English and Greek and its theoretical significance}

% \chapterDOI{} %will be filled in at production

\abstract{This paper offers a comparative study of the coordinator \emph{and}
    and the comitative preposition \emph{with} in its coordinating function.
    Greek is shown to behave differently from English in this respect and this
    is accounted for in terms of labelling\is{labelling} potential of a
    syntactic/lexical object.  The more general claims are that
labelling\is{labelling} is a locus of variation and that
labelling\is{labelling} is (still) a syntax internal process.}

\maketitle

\begin{document}\glsresetall

\section{Introduction}

One of the major proposals concerning the possible \emph{loci} of syntactic
variation is the so-called Borer--Chomsky conjecture\is{Borer--Chomsky conjecture}
which \citet{Baker2008b} formulates as follows:

\begin{quote}
All \isi{parameters} of variation are attributable to differences in features of
particular items (e.g.\ the functional heads\is{functional items}) in the lexicon.
\end{quote}

In general, it is a somewhat more restricted version that is more widely
accepted, namely that syntactic variation and parametric properties are
restricted to properties of inflectional heads only.\footnote{This is more in
line with both Chomsky's and Borer's formulations.}

In this note, I would like to suggest that the potential of a category to
supply a label to a constituent that it \emph{heads} is also a property that,
though not strictly inflectional and clearly not restricted to functional
heads, is a locus of variation across languages. The empirical argument in
favour of this position comes from the behaviour of certain coordinated
structures in \ili{English} and \ili{Greek} (and to a much lesser extent
\ili{French}). It is well known that the preposition \emph{with} in
\ili{English} also functions as a coordinator. The same is true in \ili{Greek},
but coordinations with \emph{with} pattern differently in the two languages. In
a nutshell, while in \ili{English} the first conjunct must raise out of the
\emph{with} phrase, there is no such requirement in \ili{Greek}.

In this paper I consider more closely these patterns and argue that they are
better understood if we extend \citegen{Chomsky2013} proposal on structured
coordination with \emph{and} to the case of \isi{coordination} with \emph{with}
and argue, contra \citet{Kayne1994}, that \isi{movement} of the first
conjunct is driven not by Case but by the requirements of the
labelling\is{labelling} process, and more specifically the idea that while some
categories may be able to label in some languages they may not in others.
Taking Chomsky's idea that some categories may be assigned a feature [\textsc{label}]
that nothing can remove more seriously than he probably intended, we can
imagine that this feature is an integral part of lexical items. It follows that
for categories that lack that feature, the labelling algorithm cannot identify
any of their properties for externalisation and the conceptual-intentional
system.\footnote{This is an important point to which we will return in
\Cref{consequences}.}

The paper is structured as follows: in \Cref{secdata} I present the
facts of \ili{English} concerning \emph{with}-coordinations. \Cref{laco}
develops the account of \emph{with}-coordinations in \ili{English} in labelling\is{labelling}
terms. In \Cref{greda} I turn to the \ili{Greek} data and show that the
patterns follow from the simple proposal that \ili{Greek} \emph{me} (`with') is a
labelling category. I also discuss some interpretive issues relating to
distributivity. \Cref{consequences} spells out some consequences of the
analysis.

\section{Coordination: \emph{and} and \emph{with}}\label{secdata}

The following paradigm in \ili{English} is well known:

\ea\label{andC}
    \ea[]{Sue and Sy are friends}\label{andCa}
    \ex[*]{Sue is friends and Sy}\label{andCb}
    \z
\ex\label{withC}
    \ea[]{Sue is friends with Sy}\label{withCa}
    \ex[*]{Sue with Sy are friends}\label{withCb}
    \z
\z

Examples like those in~\eqref{withC} are found with a variety of symmetric
predicates, as \citet{lakoff-peters:69} as well as \citet{dong:70} have
discussed (cf. \ref{ex:29.3}), although with varying degrees of
acceptability.

\ea\label{ex:29.3}
    \ea[]{Sue is co-workers with Sy}
    \ex[]{Sy is mates with Sue}
    \ex[]{Sue is school/bandmates with Sy}
    \ex[?]{Sy is siblings with Sue}
    \ex[]{Sue is twins with Sy}
    \ex[]{Sy is co-authors with Sue}
    \z
\z

Compare now \eqref{ex:29.3} with its version where \emph{with} is replaced by
\emph{and}.

\ea\label{morewithand}
    \ea Sue and Sy are co-workers
    \ex Sue and Sy are mates
    \ex Sue and Sy are school/bandmates
    \ex Sue and Sy are siblings
    \ex Sue and Sy are twins
    \ex Sy and Sue are co-authors
    \z
\z

The main difference between the paradigm in \eqref{ex:29.3} and that in
\eqref{morewithand} is that with \emph{and}-coordinations the whole constituent
remains together while with \emph{with} the first conjunct must move out.

Beyond nominal predicates, as above, the pattern extends to verbal symmetric
predicates such as \emph{collide} or \emph{fuck}:

\ea\label{ex:29.5}
    \ea[]{Rosetta collided with comet 67P}
    \ex[]{Rosetta and comet 67P collided}
    \ex[*]{Rosetta with comet 67P collided}
    \ex[*]{Rosetta collided and comet 67P}
    \ex[]{Sue fucks with Sy every Wednesday evening}
    \ex[*]{Sue with Sy fuck every Wednesday evening}
    \ex[]{Sue and Sy fuck every Wednesday evening}
    \ex[*]{Sue fucks and Sy every Wednesday evening}
    \z
\z

\cite{lakoff-peters:69} suggested first that the preposition \emph{with} was
functioning here as a coordinator and, moreover, the \emph{and}- and
\emph{with}-coordinations were related and should be transformationally linked
through a process of replacing \emph{and} by \emph{with} and extraposing
\emph{with NP}.  The issue of the relatedness of the two constructions as well
as the basis for \citegen{lakoff-peters:69} account was revisited, in light of
the \glsunset{LCA}\gls{LCA},\is{linear correspondence axiom} by
\cite[§6.3]{Kayne1994}, who proposed that the reason for the commonalities
between~\eqref{andCa} and~\eqref{withCa} is that they both derive from the same
underlying structure, namely \eqref{ex:29.6}.

\ea\label{ex:29.6}
    {}[DP1 [[and/with] DP2]]
\z

What sets the two constructions apart, for Kayne, is that there is a
requirement for the first conjunct to move out of the conjoined phrase in
\eqref{withCa} because it cannot be adequately Case licensed in situ. More
specifically, while a phrase coordinated with \emph{and} allows both
conjuncts to be Case licensed by virtue of the fact that the whole coordinated
constituent is in a Case-licensing position, this is not true of coordinated
phrases with \emph{with}. A somewhat different way of putting this
restriction is that, from a Case theoretic point of view, DP \isi{coordination} is
only licit if Case can be distributed to both conjuncts. In the case of
\emph{and} this appears to be so. In the case of \emph{with}, however,
this does not happen because the second conjunct is case licensed by
\emph{with} while the first one has to get Case from an external source.

The latter way of putting the relevant constraints can be made to work further,
in the sense that a constituent of the type \emph{A and B} does distribute like
its conjuncts whereas a constituent like \emph{A with B} does not. But again,
if we assume that the construction is headed by the coordinator, we would have
to suggest that in the case of \emph{with} it is still a Case assigning
preposition rather than a coordinator, which in turn casts doubt on the
analysis of these two constructions as deriving from identical underlying
structures. Moreover, under this analysis it is not clear why with different
predicates it is impossible to extract the first conjunct of a \emph{with}
coordination:

\ea[*]{Sue is {French} with Sy.}\label{sfrenchs}
\z

 For this, Kayne suggests that in order to obtain a distributive reading a coordinated phrase must be preceded by a distributor which may be overt or covert. This distributor, noted \textsc{both} following Kayne's convention, forces the distributive reading on the coordinated phrase, which is, of course equivalent to a sentential \isi{coordination}.

 \ea\label{ex:29.8}
    \textsc{both} [John and Mary] love cats $\rightarrow$ John loves cats \textsc{and} Mary loves cats.
\z

 And, of course, these cases are also fine with an overt distributor:

\ea\label{ex:29.9}
    Both John and Mary love cats.
\z

In the case of \emph{with}-coordinations, however, the distributor induces a
barrier to the \isi{movement} of the first conjunct. Thus, sentences with the
following representation are out.

\ea\label{kex}\parencite[66, example 56]{Kayne1994}\\
    John\tss{i} is human beings [\textsc{both} [[e\tss{i}] with Bill]]
\z

But it is unclear why this should be so. After all \emph{both}, as a floating
quantifier, does not induce a barrier to the \isi{movement} of its complement
\citep[cf.][]{sportiche:88}. Equally, a modifying adjunct usually does not
induce a barrier to \isi{movement} of the specifier of the category to which it
attaches. I will set aside the issues relating to interpretation and
distributivity and revisit them briefly in \Cref{distro}.

As we can see, Kayne's analysis is problematic in various respects, and yet, it
remains both plausible and attractive. In the following sections I will claim
that the basic insights can be maintained and find more elegant and general
expression in terms of the labelling\is{labelling} requirements and possibilities in these
structures.

\section{Labelling and coordination} \label{laco}

\citet{Chomsky2013} puts forward a particular proposal regarding structured
coordination (with \emph{and}), according to which coordinate structures start
as \eqref{ex:29.11}:

\ea\label{ex:29.11}
    {}[\tss{{α}} and [\tss{{β}} DP\tss{1} DP\tss{2}]]
\z

As β cannot be labelled because configurations of the type [XP YP] are
problematic for the labelling\is{labelling} algorithm (both heads are equally prominent), one
of DP\tss{1} or DP\tss{2} must raise (say DP\tss{1}) and β receives the label
of DP\tss{2}. Importantly, however, α receives the label of DP\tss{1},
reflecting the fact that the distribution of these coordinated structures is
determined by the shared label of the two coordinated elements. As Chomsky
notes, though, the construction remains headed by the conjunction which remains
visible in order to determine the structure but is not available as a label.
This entails that the whole constituent can be the target for \isi{movement} yielding
\eqref{ex:29.12} as an instance of DP \isi{movement}:\footnote{To be sure, there are various
questions surrounding Chomsky's proposal on \isi{coordination}. For example, it is
unclear what it means for the construction to headed by the coordinator, which
determines structure but does not supply a label. This requires further
clarification on the assumption that the labelling\is{labelling} algorithm identifies heads.
We set this aside for now.}

\ea\label{ex:29.12}
    {}[\tss{{DP}} Peter and Susan] are [\sout{\tss{{DP}} Peter and Susan}] teachers
\z

Assuming this to be on the right track, let us turn to the case of
\emph{with}-coor\-di\-na\-tions. Given that \eqref{ex:29.13}, modelled on
\eqref{ex:29.12} is ungrammatical, it is clear that this proposal will not
be applicable to \emph{with}-coor\-di\-na\-tions.

\ea[*]{%
    {}[\tss{{DP}} Peter with Susan] are [\sout{\tss{{DP}} Peter with Susan}]
    teachers}\label{ex:29.13}
\z

In these cases the distribution of the coordinate structure does not reflect
the distribution of their shared label (DP); in fact, it does not constitute a
well-formed constituent at all, as the data show. It follows that the derivation
will also be somewhat different. Keeping, however, as close as possible to the
proposal on \emph{and} will allow us to pinpoint the difference. The following
is a reasonable approximation of their derivation that preserves full
parallelism between the \emph{and} and the \emph{with} case. Let us assume that
DP\tss{1} and DP\tss{2} merge again like before yielding an unlabellable [XP YP]
structure. Next, \emph{with} merges with that syntactic object just like in the
case of \emph{and}. The difference, I claim, is that unlike \emph{and},
\emph{with} can provide a label for the resulting object, and we have the
following configuration:

\ea\label{ex:29.14}
    {}[\tss{{withP}} with [\tss{α} DP\tss{1} DP$2$]]
\z

At this point, DP\tss{1} must raise so that α receives the label of DP\tss{2},
yielding \eqref{ex:29.15}:

\ea\label{ex:29.15}
    {}[\tss{β} DP\tss{1} [\tss{{withP}} with [\tss{DP\tss{2}} \sout{DP\tss{1}} DP$2$]]]]
\z

Of course, the question that arises now is what label will β receive. As the
two elements of β are [DP\tss{1} withP] we are in the same situation as before
where we have a [XP YP] configuration and one of the two elements must raise.
DP\tss{1} does and following merging of further material we obtain the initial
contrast repeated here:

\ea\label{ex:29.16}
    \ea Sue is friends with Sy
    \ex Sue and Sy are friends
    \z
\z

If this is correct it is not Case but the requirement for the whole constituent
to be labelled that is responsible for the \isi{movement} of the first
conjunct. The lack of label also accounts for the fact that the whole
constituent cannot be targeted for \isi{movement}, yielding the
ungrammaticality of~\eqref{withCb}. Whether the constituent remains unlabelled
is an important question that we will pick up in \Cref{consequences}.

Although this analysis provides an account of the basic patterns, the
ungrammaticality of~\eqref{sfrenchs} remains problematic. Within the analysis
presented here, a co\-vert distributor will not do the job -- both because
assuming that it induces a barrier to \isi{movement} is not an idea that is
easy to implement in the general framework I am assuming, but also because, in
fact, even in cases like~\eqref{withCa} the reading is \emph{distributive} in
the sense that the following is a contradiction:

\ea[\#]{%
    Sue is friends with Sy but Sy is not friends with Sue.}\label{ex:29.17}
\z

With a predicate like \emph{being \ili{French}}, however, this reading is not
possible. Furthermore, the distributive reading is not really what matters, but
rather the symmetric/reciprocal one. Thus, observe the following contrast:

\ea\label{ex:29.18}
    \ea[]{Both Sebastien and Julie are French}
    \ex[*]{Both Sebastien and Julie are friends}
    \z
\z

With verbal predicates the contrast is perhaps even more telling:

\ea\label{ex:29.19}
    \ea Both Sue and Sy fucked (every/on Wednesday evening)
    \ex Both Rosetta and Galileo collided *(with comet 67P)
    \z
\z

Clearly what is missing in the meanings of the examples above is this
reciprocal/symmetrical meaning. There is no suggestion that Sue and Sy fucked
(with) each other or that Rosetta and Galileo collided with each other. Of
course, with an overt reciprocal the sentences are perfect:

\ea\label{ex:29.20}
    \ea Sue and Sy fucked each other
    \ex Rosetta and Galileo collided with each other
    \z
\z

The sentences become significantly degraded by the addition of an overt distributor:

\ea\label{ex:29.21}
    \ea[\textsuperscript{???/}*]{Both Sue and Sy fucked (with) each other}
    \ex[\textsuperscript{???/}*]{Both Rosetta and Galileo collided with each other}
    \z
\z

One way to extend the account presented here is to focus on the fact that while
\emph{and} and \emph{with} appear to perform the same function and give rise to
the same structures, it is also not true that they are synonymous.\footnote{In
\Cref{distro} I revisit this issue and propose that even if we stick
with distributivity, the results will come out right if we look more closely at
the morphology of distributivity.} Specifically, I assume that \emph{with} even
as a coordinator retains its comitative meaning and θ licenses its DP
complement (DP\tss{2} in our examples). We can then ask how is DP\tss{1}
θ-licensed.\footnote{This is a legitimate question even if we have a
coordination where we generally assume that θ licensing involves the
whole constituent. The distribution of Case inside the \emph{with}-\isi{coordination}
also does not work in the same way.} I propose here that a derivation involving
a \emph{with}-\isi{coordination} will converge only if both coordinated DPs can be
independently θ licensed.\footnote{Again, in parallel with Case.}  This
means that they will work only with two-place predicates, either verbal (like
\emph{collide}, \emph{fuck}, \emph{dance}), in which case the DP will receive a
thematic role in the subject position, or with symmetric relational nouns like
\emph{friends, co-workers} and so on where the thematic role will be available
in the nominal extended projection.\footnote{The actual mechanism is not
relevant here.} The idea, therefore, is that, unless the DP that moves out in
order to allow the [DP withP] constituent to be labelled can be thematically
licensed in its derived position, the sentence will be ungrammatical, not as a
result of lack of Case (Case can be assigned) or of lack of label, but as a
violation of the θ-criterion. Labelling is important, however, as it is
the label that allows thematic licensing in the case of
\emph{and}-coordinations and prevents it in the cases of \emph{with}, with the
results that we saw earlier. As noted earlier, there is lexical variation in
the range of elements that allow the patterns involving
\emph{with}-\isi{coordination}. So, while with a relational, symmetric noun
like \emph{friends} it works fine, with others speakers find it less acceptable
at first. Interestingly, with a noun like \emph{enemy} which allows for a
non-symmetrical reading the \emph{with} coordination is possible only in the
symmetrical reading:\footnote{Example \REF{enem} is taken from \url{http://www.davidagler.com/teaching/criticalthinking/handouts/Handout3_AdHominemFallacy.pdf}.}

\ea\label{enem}
    She is mortal enemies with John
\z

Assuming now this analysis, I turn to the corresponding \ili{Greek} facts.

\section{Greek} \label{greda}

\emph{And}-coordinations in \ili{Greek} show a behaviour similar to that of
their \ili{English} counterparts in the relevant respects, witness
(\ref{ex:29.23}--\ref{ex:29.24}):

\ea\label{ex:29.23} \ili{Greek}\\
    \gll    O Kiriakos ke o Aris ine fili.\\
            The Kiriakos and the Aris are friends \\ \label{g5}
    \glt    \enquote*{Kiriakos and Aris are friends.}
\ex\label{ex:29.24} \ili{Greek}\\
    \gll    \llap{*}O Kiriakos ine fili ke o Aris.\\
            The Kiriakos are friends and the Aris\\
    \glt    \enquote*{Kiriakos and Aris are friends.}
\z

\ili{Greek} \emph{me} \enquote*{with} also functions as a coordinator, as
in~(\ref{g1}--\ref{g2}):

\ea\label{ex:29.25} \ili{Greek}\\
    \gll    O Kostas me ton Ari ine fili.\\
            The Kostas with the Ari are friends\\ \label{g1}
    \glt    \enquote*{Kostas and Aris are friends.}
\ex\label{ex:29.26} \ili{Greek}\\
    \gll    O kostas ine filos me ton Ari.\\
            The Kostas is friend with the Aris\\ \label{g2}
    \glt    \enquote*{Kostas is friends with Aris.}
\z

At first sight, taking \ili{Greek} and \ili{English} to be basically the same,
it looks like in \ili{Greek} the first conjunct may remain in situ. From
a Case theoretic perspective this is somewhat problematic. One would wonder why
the same mechanism is not available in \ili{English}. One approach could
suggest that while we may unify \ili{Greek} and \ili{English} in terms of Case
assignment in these constructions, the \glsunset{EPP}\gls{EPP}\is{extended
projection principle} requirement of C--T must be satisfied by DP
\isi{movement} in \ili{English} while in \ili{Greek} V-to-T suffices. This is a
reasonable approach but raises the question why is it impossible to raise the
whole withP to [spec T]. The labelling\is{labelling} account developed here
provides an explanation for that. However, this question may be moot, at least
in part, given the evidence on \isi{agreement} to which we now turn. There are
some differences between \emph{with} and \emph{me}. Consider the following:

\ea\label{ex:29.27} \ili{Greek}\\
    \gll    \llap{*}O Kostas ine fili me ton Ari.\\
            The Kostas is friends with the Aris\\ \label{g3}
    \glt    \enquote*{Kostas is friends with Aris.}
\ex\label{ex:29.28} \ili{Greek}\\
    \gll    \llap{*}Ego ime fili me ton Ari.\\
            I am friends with the Ari\\ \label{g4}
    \glt    \enquote*{I am friends with Aris.}
\z

The \isi{agreement} contrast between~\eqref{g1} and~\eqref{g2} on the one hand
and~\eqref{g3} and~\eqref{g4} on the other is interesting when compared to the
agreement found in the \ili{English} \emph{friends with} construction. In the
\ili{Greek} case, plural agreement\is{agreement!number agreement} on the
predicate nominal is only triggered when the first conjunct of the [A with B]
element stays in situ. If, however, the first conjunct raises to [Spec T], then
\isi{agreement} is in the singular both on the copula in T and the predicate
nominal. Compare this to the \ili{English} \emph{friends with}
construction~\eqref{withCa} where the predicate nominal shows plural
agreement\is{agreement!number agreement} but T bears singular features (from
\isi{agreement} with the subject).  Now, given that the plural on the predicate
nominal is pretty much the only tangible evidence we can lay our hands on in
favour of the idea that the underlying structure involves a \isi{coordination},
we can take the absence of plural agreement\is{agreement!number agreement}
(together with the absence of any other factor that blocks plural
agreement\is{agreement!number agreement}) as evidence that there is no
underlying \isi{coordination} in \ili{Greek}, and the right analysis
of~\eqref{g2} is roughly \eqref{ex:29.29}:

\ea\label{ex:29.29}
    \begin{tikzpicture}[baseline=(root.base)]

        \Tree   [.\node(root){TP};
                    [.DP \edge[roof]; {O Kostas} ]
                    [.T
                        [.T ine ]
                        [.SC
                            [.AP filos ]
                            [.PP
                                me
                                [.DP \edge[roof]; {to Giani} ]
                            ]
                        ]
                    ]
                ]

    \end{tikzpicture}
\z

\emph{Friends with} construction is not available in \ili{Greek}. Under a Case
theoretic approach, this is problematic given that \emph{me} assigns Case to
its complement DP while DP\tss{1} has its Case valued externally. So even
pursuing that path one would have to find out why \ili{Greek} allows this type
of Case valuation in cases that look otherwise equivalent.

Given the discussion above and the \isi{agreement} facts, it is, I suggest,
reasonable to propose that the difference between \ili{Greek} and \ili{English} regarding
\emph{with}-coordinations should be located in the labelling\is{labelling} potential of
\emph{with/me}.

In the previous section we saw that in \ili{English} \emph{with} was different
from \emph{and} in that it could supply a label. I want now to propose that in
\ili{Greek} \emph{me} is exactly the same as \emph{ke} \enquote*{and} in terms
of labelling\is{labelling} potential,\footnote{They are different in other
    ways, see \Cref{distro}.} i.e. neither can supply a label (in other words
    nether carries the feature [\textsc{label}]), and, as a result, it is not surprising
    that the behaviour of \emph{me}-coordinations in \ili{Greek} is similar to
    that of \emph{and}-coordinations\is{coordination} (in \ili{Greek} and
    \ili{English}).  Assuming this, the patterns follow.

Consider first the fact that the whole constituent will be labelled DP and as a
result can be targeted for \glsunset{EPP}\gls{EPP}\is{extended projection
principle} driven \isi{movement} and for Case valuation.  Concerning Case, as
we saw above, \emph{me} will Case license DP\tss{2} while DP\tss{1} will have
its Case valued via \isi{Agree} with T. The following examples show that the
whole DP can appear preverbally in subject position with different nominal or
prepositional predicates:

\ea\label{ex:29.30} \ili{Greek}\\
    \gll    Ego me ton patera mu imaste sinehia se sigrusi.\\
            I with the father mine are always in collision\\
    \glt    \enquote*{I am always fighting with my father.}
\ex\label{ex:29.31} \ili{Greek}\\
	\gll    Ego me ton Kosta imaste aderfia.\\
            I with the Kostas are siblings\\
    \glt    \enquote*{Kostas and I are siblings.}
\ex\label{ex:29.32} \ili{Greek}\\
	\gll    Ego me ton Apostoli imaste panda antipali.\\
            I with the Apostolis are always rivals\\
    \glt    \enquote*{Apostolis and I are always rivals.}
\z

Assuming further that in some way coordinated phrases are marked as formally
plural, \isi{agreement} both with the predicate nominal and T is expected to be
in the plural. This prediction is borne out.

Furthermore, we predict that these coordinated structures will be available
with a wide variety of verbal predicates too; in other words, not just with the
symmetric ones with which they co-occur in \ili{English}. Again the prediction
is borne out as the following examples show:\footnote{The examples
(\ref{g7}--\ref{g12}) were found with a simple Google search.}

\ea\label{ex:29.33} \ili{Greek}\\
	\gll    O tragudistis me ti sizigo tu tu ehun megali adinamia.\\
            The singer with the spouse his to-him have great weakness\\ \label{g7}
    \glt    \enquote*{The singer and his wife have a weak spot for him.}
\ex\label{ex:29.34} \ili{Greek}\\
	\gll    O Kostas me ti Marina, pu ehun molis padrefti, benun mesa sto
    saloni.\\
            The Kostas with the Marina, who have just married, enter in the living-room\\ \label{g8}
    \glt    \enquote*{Kostas and Marina, who just got married, enter the living
    room.}
\ex\label{ex:29.35} \ili{Greek}\\
	\gll    O Nikos me ti Maria ehun dio pedia.\\
            The Nikos with the Maria have two children\\ \label{g9}
    \glt    \enquote*{Nikos and Maria have two children.}
\ex\label{ex:29.36} \ili{Greek}\\
	\gll    O Sakis me ti Frini apoktisan pedi.\\
            The Sakis with the Frini obtained child\\ \label{g10}
    \glt    \enquote*{Sakis and Frini had a child.}
\ex\label{ex:29.37} \ili{Greek}\\
	\gll    O Panagiotis me ti Hrisa ehun anagagi to kreopolio tus se horo
    sinathrisis.\\
            The Panagiotis with the Hrisa have elevated the butcher's theirs to space rally\\ \label{g11}
    \glt    \enquote*{Panagiotis and Hrisa have turned their butcher's shop to
    a major gathering place.}
\ex\label{ex:29.38} \ili{Greek}\\
	\gll    O Grigoris me ton Petro kserun pos tha se odigisoun.\\
            The Grigoris with the Petros know how will you drive\\ \label{g12}
    \glt    \enquote*{Grigoris and Petros know how to drive you around.}
\ex\label{ex:29.39} \ili{Greek}\\
	\gll    Telika i Rihana me to Saudarava ine mazi edo ke mines.\\
            Finally the Rihana with the Saudi are together here and months\\ \label{g13}
    \glt    \enquote*{In the end Rihana and the Saudi man have been together
    for months.}
\z

The interpretation of these examples is dependent on the predicate; if the
predicate allows for a symmetric reading like~\eqref{g8}, where if A is married
to B then B is also married to A, then this is what we obtain. If the predicate
allows or requires a group reading, like (\ref{g11}--\ref{g12}), this what we
get. And finally, if the predicate allows or requires a distributive reading,
like~\eqref{g7} or one reading of~\eqref{g9} this is again what we have.

Under the simple proposal that \emph{me} is a non-labelling head the data above
are all expected. Let me now turn to a somewhat complicating factor, namely
distributivity.

\subsection{A complication: Distributivity} \label{distro}

There seems to be one significant difference between \emph{ke} and \emph{me} in
\ili{Greek}. It is well known that in \ili{Greek}, like in \ili{French}, the
coordinator can appear in front of both coordinated constituents:

\ea\label{ex:29.40} French\\
    \gll    Pierre conna\^it et Isabelle et Marie.\\
            Pierre knows and Isabelle and Marie\\
    \glt    \enquote*{Pierre knows both Isabelle and Marie.}
\pagebreak\ex\label{ex:29.41} \ili{Greek}\\
	\gll    O Kostas gnorizi ke ti Maria ke tin Eleni.\\
            The Kostas knows and the Maria and the Eleni\\
    \glt    \enquote*{Kostas knows both Maria and Eleni.}
\z\largerpage

\citet[146, fn.\ 16]{Kayne1994} for \ili{French} and \citet{cms:15} for \ili{Greek} have
argued that the initial (outer) occurrence of the coordinator is in fact a
distributive operator. Although this is generally true in the sense that the
initial \emph{ke/et} yields a distributive reading it is also true that this is
only the case when the second (inner) coordinator is \emph{and/ke/et}. Thus,
in \ili{Greek}, with a \emph{me}-\isi{coordination} no distributive readings are induced by
the presence of an initial \emph{ke}, compare:

\ea\label{ex:29.42} \ili{Greek}\\
	\gll    Ke o Sakis ke i Sula sikosan ena trapezi.\\
            And the Sakis and the Sula lifted a table\\
    \glt    \enquote*{Both Sakis and Sula lifted a table.} \hfill \textsc{distributive}
\ex\label{ex:29.43} \ili{Greek}\\
	\gll    Ke o Sakis me ti Sula sikosan ena trapezi.\\
            And the Sakis with the Sula sikosan ena trapezi\\
    \glt    \enquote*{Sakis and Sula lifted a table.} \hfill \textsc{collective}
\z

Now perhaps it is the comitative meaning of \emph{me} (which was suggested in
\Cref{laco} for \ili{English} and is presumably also valid for \ili{Greek}) that somehow
blocks the distributive reading. One way of putting this is to suggest that,
semantically, the output of a \emph{me}-\isi{coordination} is a \emph{group}
individual, acting in part as an atom, whereas this is not necessary for
\emph{ke}-coordinations, whose semantic value may be that of a \emph{group} (in
which case there is no difference with \emph{me}) but can also be an individual
of type \textsc{sum}, which would be an appropriate argument for the distributive
operator. However, examples like \eqref{ex:29.44} seem to suggest
otherwise, in the sense that, as things stand, there is no immediate suggestion
that the two teams form a group in a relevant sense:\footnote{This is perhaps
    too strong.  The two teams might form a group in the sense that they are
    the two \ili{Greek} teams in the relevant international championship. I
will set this aside for this paper.}

\ea\label{ex:29.44} \ili{Greek}\\
	\gll    O Olimpiakos me ton Panathinaiko kserun pia apenandi se pies omades
    tha agonistun.\\
            The Olimpiakos and the Panathinaikos know at-last against to which teams will play\\
    \glt    \enquote*{Olimpiakos and Panathinaikos have at last found out which
    teams they will face.}
\z

The reading of \eqref{ex:29.44} is distributive in the sense that it
corresponds to a sentential conjunction \eqref{ex:29.45}:

\ea\label{ex:29.45}
    Olympiakos knows which team it will face \textsc{and} Panathinaikos knows which team
    it will face.
\z

Now adding an initial \emph{ke} to \eqref{ex:29.44} does not have the
desired effect:

\ea\label{ex:29.46} \ili{Greek}\\
	\gll    Ke o Olimpiakos me ton Panathinaiko kserun pia apenandi se pies
    omades tha agonistun.\\
            And the Olimpiakos with the Panathinaikos know at-last against to which teams will play\\
    \glt    \enquote*{Olimpiakos and Panathinaikos \textbf{\textsc{also}} have at last
    found out which teams they will face (as well as some other group of
teams).}
\z

In this case the reading is that of the additive \emph{ke}.\footnote{For more
details on the additive \emph{ke}, see \citet{cms:15} and references therein.}

Another issue with the idea that the initial \emph{ke} is the distributive
operator applying to an argument of \textsc{sum} type is that \emph{ke}, \emph{qua}
distributive operator, is not available with plurals, which are routinely
thought of as carrying the type of sums (\cite{link2} and many more after him).
Interestingly this is not true for \ili{English} \emph{both}:\footnote{In
    \ili{French} the relevant sentences are altogether ungrammatical so we will
    not pursue the comparison further although the question why the
    distributive \emph{et} cannot appear with plurals in any position is an
    intriguing one:

\begin{exe}\label{ex:29.47}
    \exi{(i)} \ili{French}\\
    \gll     \llap{*}Et les enfants ont soulev\'{e} une table\\
             and the children have lifted a table\\
    \glt     \enquote*{The children have lifted a table.} (intended: each)
%\end{exe}
%
%\begin{exe}
    \exi{(ii)} \ili{French}\label{ex:29.48}\\
    \gll     \llap{*}Jean conna\^{i}t et les enfants\\
             Jean knows and the children \\
    \glt     intended: \enquote*{Jean knows each child.}
\end{exe}}

\ea\label{ex:29.49} \ili{Greek}\\
	\gll    Ke ta pedia efagan gemista.\\
            And the children ate gemista\\ \label{g45}
    \glt    \enquote*{The children \textsc{too} ate gemista.}
\ex\label{ex:29.50}
    Both children ate gemista
\z

Again the \emph{ke} on~\eqref{g45} is the additive \emph{ke} and does not give
the desired distributive reading, unlike what we see in \eqref{ex:29.50}.

Setting aside this concern, these patterns can be understood in two ways which
probably boil down to the same insight. On the one hand, as suggested earlier,
we can think of inner \emph{and/ke/et} as sum forming operators and outer
\emph{ke/et} as distributors acting upon these sums. In contrast \emph{with/me}
are group forming operators whose outcome behaves in the relevant respects as
an atom and therefore the distributor cannot act on them in the same way. This
would mean that the reason why initial \emph{ke} followed by a \emph{with}
coordination can only be read as additive falls together with
\eqref{ex:29.51}:

\ea\label{ex:29.51} \ili{Greek}\\
	\gll    Ke i epitropi apofasise tin isvoli stin Amorgo.\\
            And the committee decided the invasion to-the Amorgos\\
    \glt    \enquote*{The committee (as well as some other organisation)
    decided the invasion of Amorgos.}
\z

The alternative way of analysing these patterns is to suggest that the
distributive operator is in fact the discontinuous morpheme:

\ea\label{ex:29.52}
    \ea Both \ldots{} and
    \ex Ke \ldots{} ke
    \ex Et \ldots{} et
    \z
\z

Again this idea predicts that adding \emph{both} or \emph{ke} in front of a
\emph{with/me}-\isi{coor\-di\-na\-tion} will not yield a distributive reading simply
because, at least in these cases, it is just not the right morpheme for the
intended meaning. I think that in this way the ungrammaticality of Kayne's
example~\eqref{kex}, repeated here, is explained too:

\ea\label{ex:29.53}
    John\tss{i} is human beings [\textsc{both} [[e\tss{i}] with Bill]]
\z

While Kayne is right that distributivity is the key to understanding the
judgement, it is not because a covert \textsc{both} blocks the extraction. Rather, it is
because the distributive reading does not arise in these cases because the
lexical material is just not right.

\section{Some consequences}\label{consequences}

Let us take stock. I argued so far in this paper that a number of differences
in the syntax of \isi{coordination} both within and across languages can be
understood in terms of the labelling\is{labelling} potential of different
categories and the labelling algorithm. The account developed here raises a
number of questions primarily about the role of labels in syntactic
derivations.

A particular point of debate regarding labelling\is{labelling} going back to
the early days of minimalism is whether labels are mere tags onto pieces of
structure serving to identify them as a potential targets for operations such
as internal \isi{Merge} or \textsc{agree} at least,\footnote{The question of
external merge is also relevant in terms of the elements that are identified
for \isi{Merge}.} or active drivers of the derivation.
\citet{Chomsky1993,Chomsky1995} took the former view.  A different view was
taken by \cite{adger-tsoulas:99}, who proposed that labels are complex and
include category determining features from both merged elements, i.e. Merge$(\text{α},
\text{β}) \rightarrow [_{\{\text{α,β}\}}\ \text{α}, \text{β} ]$. Crucially, the label $ \{\text{α,β}\}$ was taken
to be semi-uninterpretable in the sense that one of the two categorial features
that make it up (α and β) had to be eliminated.  Eliminating that feature was
done in the standard way, by seeking a goal in the numeration or the sub-array,
agreeing, and merging it with the existing structure or, by internal merge,
raising an element with the required specification. In that proposal,
computation was driven by the labels, whether on heads or intermediate
projections. Although Chomsky's recent proposals on labelling and the one from
\citet{adger-tsoulas:99} differ in many respects, they converge on the idea
that determining the label of a particular part of the structure is a driving
force for computation and that in principle labelling need not obey
endocentricity.  They diverge on two important conceptual points, namely (a)
whether the output of merge needs to be always labelled, and (b) what are
labels required for. Regarding the former, \citet[6]{Chomsky2015} is
particularly clear on this point:\glsunset{LA}

\begin{quote}

Crucially, \gls{LA} does not yield a new category as has been assumed in PSG
and its various descendants, including X$'$ theory. Under LA, there is no
structure [\tss{α} X], where α is the label of X. \gls{LA} simply determines a
property of X for externalization and CI. It is therefore advisable to abandon
the familiar tree notations, which are now misleading.  Thus in the description
of an [XP, [YP, ZP]] structure, there is no node above either of the two merged
constituents. There is no label for the root of the branching nodes.

\end{quote}

Taking this at face value, it means that not every output of merge operations
will be labelled. A question we might ask about this approach is what happens
to elements such as [ α, β ] when \gls{LA} has not identified a property for
externalisation and CI. The issue is puzzling. Imagine that there is some
element X for which the Labelling algorithm as identified no property (I
suppose that this would be its label) for externalisation and CI. What would
that actually mean? In terms of externalisation it would mean that the element
would not be pronounced. This is the reasonable understanding of the idea (from
\citet{Chomsky2015} that copies do not label. In other words the algorithm will
identify no property of copies relevant to externalisation. \emph{Wanna}
contraction aside, this seems correct. But what of CI? Would one expect that
such an element would be invisible also to the interpretive mechanisms? This
seems problematic. Focusing on the cases of interest in this paper, both
\emph{and/ke}- and (in \ili{Greek} at least) \emph{me}-coordinations would be such
that the coordinator would provide no relevant property for externalisation and
CI. If the reasoning based on copies is on the right track, then the
non-labelling nature of the coordinators is a clear counterexample (they are
after all externalised). But setting externalisation aside, in the case of CI
it is unclear, in this case, how a structure [DP\tss{1} and DP\tss{2}] would be
interpreted. What does seem clear is that it is a property of the conjunction
that is preeminent in the interpretation, namely whatever it is that turns that
constituent into a plural (sum) entity. Assume for concreteness that the
semantics for DP conjunction corresponds to set formation, or more precisely
set-product formation, defined in its general form as follows
\citep[241]{heycock-zamparelli:05a}:

\ea\label{ex:29.53b}
    Set product \textsc{(sp)}\\
    \textsc{sp}$(S^1,\dots, S^n) =_{\text{def}} \{X:X=A^1\cup \dots \cup A^n, A^1 \in S^1,\dots, A^n \in S^n\}$
\z

The way this works is by taking one element from the denotation of each of the
two conjoined elements and yielding their union for all elements of these sets.
This is the property that is relevant to CI, rather than the DP label that, as
we saw, is assigned by the labelling\is{labelling} algorithm. The DP label (or
at the very least the lack of label deriving from the conjunction), however, is
precisely what accounts for the syntactic patterns. Thus, if the reasoning is
correct, we are led to rethink the labelling\is{labelling} process as follows:
labels in part drive syntactic computation but in crucial respects do not
represent properties for CI and externalisation. There is a mismatch between
the label relevant to the derivation itself and the CI/semantically relevant
one. Labels are necessary and the labelling\is{labelling} algorithm is a tool
that affords insightful understandings of syntactic patterns, but labels do not
determine interface interpretation and do not reflect interface properties.
Often in fact, as in the cases analysed in this paper, the syntactic label is
at odds with the semantically relevant one.

\section{Conclusion}

In this paper I tried to rethink the properties of two types of
\isi{coordination} in \ili{English} and \ili{Greek}. I argued that the
different behaviour of \emph{and} and \emph{with}-coor\-di\-na\-tion in
\ili{English} are the result of the fact that while \emph{and} does not provide
a syntactic label \emph{with} does. In \ili{Greek}, however, neither did,
resulting in different behaviours. If I am correct we probably also have to
accept two higher level conclusions. First, that the (non)-labelling nature of
a category can capture linguistic variation and perhaps is a parametric
property. Given that this is not an inflectional category, if I am correct,
then there is evidence for variation that, although ultimately located in the
lexicon if we assume that there is a feature [\textsc{label}], concerns the only thing
that is determined internally to the computational system. The second
conclusion, connected directly to the first, is that labelling is a process
necessary for the syntactic computation and is neither determined by nor
determines interface properties.

\printchapterglossary{}

\section*{Acknowledgements}

A number of people have very patiently discussed with me the material in this
paper, have shared very generously their judgements, and have occasionally
stopped me from making some important mistakes. In alphabetical order, I want
to thank: Kook-Hee Gil, Nino Grillo, Ekali Kostopoulos, Margarita Makri,
Dimitris Michelioudakis, Gillian Ramchand, Peter Sells, Hanna de Vries, Rebecca
Woods, and Norman Yeo. Unless the mistakes they prevented were not mistakes
then they are not responsible for any shortcomings. Theresa Biberauer’s
combination of encouragement and understanding has been more instrumental to
the completion of this paper than anything else (but she should not be blamed
for it). I am extremely happy to offer this to Ian on his birthday and raise
\emph{n} glasses to many happy rethinks.

\newpage\largerpage
{\sloppy\printbibliography[heading=subbibliography,notkeyword=this]}

\end{document}
